% 理想二维孔隙 NMR T2：精确解析解与快扩散近似（含实际模拟参数）
% 编译：xelatex nmr_analytic_solution.tex（两次）
\documentclass[11pt,a4paper]{article}
% ctex/zhnumber not installed locally; xeCJK + Windows fonts give the same Chinese typesetting.
\usepackage{xeCJK}
\setCJKmainfont[BoldFont=SimHei,ItalicFont=KaiTi]{Microsoft YaHei}
\setCJKsansfont{Microsoft YaHei}
\renewcommand{\tablename}{表}
\renewcommand{\refname}{参考文献}
\renewcommand{\contentsname}{目录}
\usepackage[margin=2.2cm]{geometry}
\usepackage{amsmath,amssymb,bm}
\usepackage{siunitx}
\usepackage{booktabs}
\usepackage{array}
\usepackage[round,semicolon]{natbib}
\setcitestyle{aysep={,}}
\usepackage[colorlinks=true,linkcolor=blue,citecolor=blue,urlcolor=blue]{hyperref}
\sisetup{per-mode=symbol,group-digits=false}
\newcommand{\TB}{T_{2B}}
\newcommand{\Ts}{T_{2S}}

\title{理想二维孔隙中 NMR $T_2$ 的精确解析解与快扩散近似\\[2pt]
\large ——以“大孔–狭长喉道–小孔”模拟案例的参数代入为例}
\author{}
\date{}

\urlstyle{same}
\begin{document}
\maketitle

\section{物理模型与参数}

本节对应数值模拟中的孤立孔参照（isolated-pore reference）：大孔、小孔分别为二维圆盘（半径 $R$），
喉道为宽 $w=2h$ 的二维狭缝（平板，长度远大于宽度）。模型为二维（每单位纵深），即沿 $z$ 方向无限平移不变的挤出体，
因此圆盘对应无限长圆柱截面，狭缝对应无限大平行平板。

\begin{table}[h]
\centering
\caption{模拟输入参数（与 \texttt{Auto\_NMR.py} 默认值一致，\citealp{Zhou2025}）}
\begin{tabular}{llll}
\toprule
符号 & 含义 & 取值 & 计算单位 \\
\midrule
$D$ & 水的自扩散系数 & \SI{2.0e-9}{m^2/s} & \SI{2.0}{\micro m^2/ms} \\
$\rho_2$ & 横向表面弛豫率 & \SI{5}{\micro m/s} & \SI{0.005}{\micro m/ms} \\
$\TB$ & 体弛豫时间 & \SI{3}{s} & \SI{3000}{ms} \\
$R_B$ & 大孔半径 & \SI{20}{\micro m} & \\
$R_S$ & 小孔半径 & \SI{5}{\micro m} & \\
$h$ & 喉道半宽（$w=\SI{0.1}{\micro m}$） & \SI{0.05}{\micro m} & \\
\bottomrule
\end{tabular}
\end{table}

\section{控制方程与边界条件}

无外加梯度时，横向磁化 $M(\bm r,t)$ 满足 Bloch–Torrey 扩散–弛豫方程\citep{Torrey1956}
\begin{equation}
\frac{\partial M}{\partial t}=D\nabla^{2}M-\frac{M}{\TB},\qquad \bm r\in\Omega ,
\label{eq:bt}
\end{equation}
孔壁 $\partial\Omega$ 上为表面弛豫的 Robin 边界条件\citep{Senturia1970,BrownsteinTarr1979}
\begin{equation}
D\,\bm n\cdot\nabla M+\rho_2 M=0,\qquad \bm r\in\partial\Omega ,
\label{eq:robin}
\end{equation}
其中 $\bm n$ 为外法向。初始条件为均匀磁化 $M(\bm r,0)=M_0$；测得信号为孔内积分
\begin{equation}
S(t)=\int_\Omega M(\bm r,t)\,\mathrm d\Omega .
\end{equation}

\section{分离变量与通解}

令 $M(\bm r,t)=e^{-t/\TB}\,m(\bm r,t)$，则体弛豫完全分离，$m$ 满足纯扩散方程 $\partial_t m=D\nabla^2 m$
及同样的 Robin 条件 \eqref{eq:robin}。再令 $m=\phi(\bm r)\,e^{-t/\Ts}$，得到本征值问题
\begin{equation}
\nabla^{2}\phi_n+\frac{\xi_n^{2}}{a^{2}}\phi_n=0,\qquad
D\,\bm n\cdot\nabla\phi_n+\rho_2\phi_n=0,
\label{eq:eig}
\end{equation}
其中 $a$ 为特征长度（圆盘取 $a=R$，平板取 $a=h$）。定义无量纲表面弛豫数（Brownstein–Tarr 数）
\begin{equation}
\kappa=\frac{\rho_2\,a}{D}.
\end{equation}
第 $n$ 个模态的弛豫速率与归一化信号为\citep{BrownsteinTarr1979,Grebenkov2007}
\begin{equation}
\frac{1}{T_{2,n}}=\frac{1}{\TB}+\frac{D\,\xi_n^{2}}{a^{2}},\qquad
\frac{S(t)}{S(0)}=\sum_{n=0}^{\infty}A_n\exp\!\left(-\frac{t}{T_{2,n}}\right),\qquad \sum_n A_n=1 .
\label{eq:modes}
\end{equation}

\subsection{二维圆盘（孔体）}

轴对称本征函数为 $\phi_n(r)=J_0(\xi_n r/R)$。由 $J_0'(x)=-J_1(x)$，Robin 条件在 $r=R$ 处给出特征方程
\begin{equation}
\boxed{\;\xi_n\,\frac{J_1(\xi_n)}{J_0(\xi_n)}=\kappa=\frac{\rho_2 R}{D}\;}
\label{eq:disk}
\end{equation}
各模态的信号权重为\citep{BrownsteinTarr1979,CarslawJaeger1959}
\begin{equation}
A_n=\frac{4J_1^{2}(\xi_n)}{\xi_n^{2}\left[J_0^{2}(\xi_n)+J_1^{2}(\xi_n)\right]}
=\frac{4\kappa^{2}}{\xi_n^{2}\left(\xi_n^{2}+\kappa^{2}\right)} .
\end{equation}

\subsection{二维狭缝（喉道，平行平板）}

以狭缝中线为 $y=0$，偶对称本征函数为 $\phi_n(y)=\cos(\xi_n y/h)$，Robin 条件在 $y=\pm h$ 处给出
\begin{equation}
\boxed{\;\xi_n\tan\xi_n=\kappa=\frac{\rho_2 h}{D}\;},\qquad
A_n=\frac{2\sin^{2}\xi_n}{\xi_n\left(\xi_n+\sin\xi_n\cos\xi_n\right)} .
\label{eq:slab}
\end{equation}

\section{快扩散极限}

当 $\kappa\ll1$ 时，扩散穿越孔隙的时间 $a^2/D$ 远短于表面弛豫时间 $a/\rho_2$，孔内磁化保持近似均匀，
信号几乎只含最低模态（$A_0\to1$），这就是 Brownstein–Tarr 快扩散区\citep{BrownsteinTarr1979}。

对圆盘，利用 $J_1(x)/J_0(x)=\tfrac{x}{2}+\tfrac{x^{3}}{16}+O(x^{5})$，式 \eqref{eq:disk} 变为
$\tfrac{\xi_0^{2}}{2}+\tfrac{\xi_0^{4}}{16}\approx\kappa$，故
\begin{equation}
\xi_0^{2}=2\kappa\left(1-\frac{\kappa}{4}\right)+O(\kappa^{3})
\;\;\Longrightarrow\;\;
\frac{1}{\Ts}=\frac{D\xi_0^{2}}{R^{2}}\approx\frac{2\rho_2}{R}\left(1-\frac{\kappa}{4}\right).
\end{equation}
对平板，利用 $x\tan x=x^{2}+\tfrac{x^{4}}{3}+O(x^{6})$，得 $\xi_0^{2}\approx\kappa(1-\kappa/3)$，
$1/\Ts\approx(\rho_2/h)(1-\kappa/3)$。保留首项即得经典的快扩散关系\citep{BrownsteinTarr1979,Kleinberg1994,Coates1999}
\begin{equation}
\boxed{\;\frac{1}{T_2}\approx\frac{1}{\TB}+\rho_2\frac{S}{V}\;},\qquad
\frac{S}{V}=
\begin{cases}
2/R, & \text{二维圆盘（无限长圆柱）},\\[2pt]
1/h, & \text{二维狭缝（平行平板）}.
\end{cases}
\label{eq:fd}
\end{equation}
反过来，由测得的 $T_2$ 可换算孔隙尺寸：
\begin{equation}
R=\frac{2\rho_2}{\dfrac{1}{T_2}-\dfrac{1}{\TB}},\qquad
h=\frac{\rho_2}{\dfrac{1}{T_2}-\dfrac{1}{\TB}} .
\label{eq:inverse}
\end{equation}
同一 $T_2$ 若按圆盘公式换算，平板（喉道）会得到 $R=2h$，即狭缝模态在“圆孔换算轴”上被放大为其半宽的两倍。

\section{代入模拟参数的完整求解过程}

\subsection{大孔 $R_B=\SI{20}{\micro m}$}

\textbf{第 1 步：判断扩散区。}
\[
\kappa=\frac{\rho_2R}{D}=\frac{0.005\times 20}{2.0}=0.05\ll 1\quad\text{（快扩散区）}.
\]
\textbf{第 2 步：数值求解特征方程 \eqref{eq:disk}。}
\[
\xi_0\,\frac{J_1(\xi_0)}{J_0(\xi_0)}=0.05
\;\Longrightarrow\;
\xi_0=0.314262,\qquad \xi_0^{2}=0.0987604 .
\]
\textbf{第 3 步：表面弛豫与总弛豫时间。}
\begin{align*}
\frac{1}{\Ts}&=\frac{D\xi_0^{2}}{R^{2}}=\frac{2.0\times0.0987604}{20^{2}}=\SI{4.93802e-4}{ms^{-1}},
&\Ts&=\SI{2025.10}{ms},\\
\frac{1}{T_2}&=\frac{1}{\Ts}+\frac{1}{\TB}=4.93802\times10^{-4}+3.33333\times10^{-4}=\SI{8.27135e-4}{ms^{-1}},
&T_2&=\SI{1208.99}{ms}.
\end{align*}
\textbf{第 4 步：模态权重（检验单指数假设）。}
\[
A_0=\frac{4\kappa^2}{\xi_0^2(\xi_0^2+\kappa^2)}=0.999949;\qquad
\xi_1=3.84473,\;\; T_{2,1}=\SI{13.47}{ms},\;\; A_1=4.6\times10^{-5}.
\]
\textbf{第 5 步：快扩散近似对照。}
\begin{align*}
\frac{1}{\Ts}&\approx\frac{2\rho_2}{R}=\frac{2\times0.005}{20}=\SI{5.0e-4}{ms^{-1}},
&\Ts&\approx\SI{2000}{ms},\\
\frac{1}{T_2}&\approx5.0\times10^{-4}+3.33333\times10^{-4}=\SI{8.33333e-4}{ms^{-1}},
&T_2&\approx\SI{1200.0}{ms}.
\end{align*}
相对精确解偏差 $-0.74\%$；加一阶修正 $(1-\kappa/4)$ 后 $\Ts\approx\SI{2025.3}{ms}$，偏差降至 $0.01\%$。

\subsection{小孔 $R_S=\SI{5}{\micro m}$}

\begin{align*}
\kappa&=\frac{0.005\times5}{2.0}=0.0125,
&\xi_0&=0.157867,\quad \xi_0^{2}=0.0249220,\\
\frac{1}{\Ts}&=\frac{2.0\times0.0249220}{5^{2}}=\SI{1.99376e-3}{ms^{-1}},
&\Ts&=\SI{501.56}{ms},\\
\frac{1}{T_2}&=\frac{1}{501.56}+\frac{1}{3000}=\SI{2.32709e-3}{ms^{-1}},
&T_2&=\SI{429.72}{ms},\quad A_0=0.999997 .
\end{align*}
快扩散近似：$1/\Ts\approx 2\times0.005/5=\SI{2.0e-3}{ms^{-1}}$，$\Ts\approx\SI{500}{ms}$，
$T_2\approx\SI{428.57}{ms}$，偏差 $-0.27\%$。

\subsection{喉道（狭缝）$h=\SI{0.05}{\micro m}$}

\begin{align*}
\kappa&=\frac{0.005\times0.05}{2.0}=1.25\times10^{-4},
&\xi_0\tan\xi_0&=\kappa\;\Rightarrow\;\xi_0=0.0111801,\\
\frac{1}{\Ts}&=\frac{D\xi_0^{2}}{h^{2}}=\frac{2.0\times1.24995\times10^{-4}}{0.05^{2}}=\SI{0.0999958}{ms^{-1}},
&\Ts&=\SI{10.0004}{ms},\quad T_2=\SI{9.9672}{ms}.
\end{align*}
快扩散近似：$1/\Ts\approx\rho_2/h=0.005/0.05=\SI{0.1}{ms^{-1}}$，$T_2\approx\SI{9.9668}{ms}$，偏差 $-0.004\%$。

\subsection{汇总与数值模拟对照}

\begin{table}[h]
\centering
\caption{孤立孔 $T_2$：精确解、快扩散近似与 pyGIMLi 数值解（单位 ms）}
\begin{tabular}{lcccccc}
\toprule
孔隙 & $\kappa=\rho_2a/D$ & $\xi_0$ & 精确 $T_2$ & 快扩散 $T_2$ & 偏差 & \shortstack{数值模拟 $T_2$\\（单指数拟合）} \\
\midrule
大孔 $R=\SI{20}{\micro m}$ & 0.05 & 0.314262 & 1208.99 & 1200.00 & $-0.74\%$ & 1209.26 \\
小孔 $R=\SI{5}{\micro m}$ & 0.0125 & 0.157867 & 429.72 & 428.57 & $-0.27\%$ & 429.83 \\
喉道 $h=\SI{0.05}{\micro m}$ & $1.25\times10^{-4}$ & 0.0111801 & 9.9672 & 9.9668 & $-0.004\%$ & 10.03 \\
\bottomrule
\end{tabular}
\end{table}

数值模拟为 pyGIMLi 三角网格上的隐式 Euler 有限元解（$\Delta t=\SI{0.25}{ms}$），对归一化衰减在 $0.9$–$0.05$ 区间做单指数拟合；
与精确解的偏差分别为 $+0.02\%$、$+0.03\%$、$+0.6\%$（喉道偏差主要来自时间步长 $\Delta t/T_2\approx2.5\%$ 的隐式 Euler 截断误差）。

\subsection{由模拟 $T_2$ 反算孔隙尺寸}

数值模拟直接输出的是回波衰减 $M(t)$，$T_2$ 需要从中提取。表 2 的“数值模拟 $T_2$”是对每个孤立孔\emph{自身}的衰减做单指数拟合
（$M/M_0$ 在 $0.9$–$0.05$ 区间对 $\ln M$ 线性拟合）得到的，未经 NNLS 反演；孤立孔主模态权重 $A_0>0.9999$，故可视为单指数。
代入式 \eqref{eq:inverse}：
\begin{align*}
R_B&=\frac{2\times0.005}{1/1209.26-1/3000}=\frac{0.01}{4.93622\times10^{-4}}=\SI{20.26}{\micro m}\quad(+1.3\%),\\
R_S&=\frac{2\times0.005}{1/429.83-1/3000}=\frac{0.01}{1.99315\times10^{-3}}=\SI{5.017}{\micro m}\quad(+0.3\%),\\
h&=\frac{0.005}{1/10.03-1/3000}=\frac{0.005}{9.94047\times10^{-2}}=\SI{0.0503}{\micro m}\quad(+0.6\%).
\end{align*}
大孔反算偏大 $1.3\%$，即快扩散近似本身的误差（$\kappa/4\approx1.25\%$）。

\subsection{$T_2$ 提取方式对反算尺寸的影响}

同一组模拟回波可以用以下四种方式（方式 1–4）提取 $T_2$。NNLS 反演均采用 \texttt{nmr\_t2} 的 L-curve 准则（SNR $=1000$，200 个分箱，$1$–$10^5$ ms），
$T_{2\mathrm{lm}}$ 取该峰（连续非零支撑区）内的对数平均：
\begin{enumerate}
\item 直接单指数拟合孤立孔自身衰减；
\item 孤立孔自身衰减单独做 NNLS 反演；
\item 将三种孤立孔的衰减按面积相加（$B+S+$ 两段喉道）后做 NNLS 反演；
\item 连通链（考虑孔间扩散耦合）的衰减做 NNLS 反演。
\end{enumerate}

\begin{table}[h]
\centering
\caption{四种 $T_2$ 提取方式的反算尺寸（单位 \si{\micro m}；括号内为相对真值偏差；NNLS 列为“主峰 / $T_{2\mathrm{lm}}$”）}
\small
\begin{tabular}{lcccc}
\toprule
孔隙（真值） & \shortstack{方式 1\\直接拟合} & \shortstack{方式 2\\单独 NNLS} & \shortstack{方式 3\\合成信号 NNLS} & \shortstack{方式 4\\连通链 NNLS} \\
\midrule
大孔（20） & \shortstack{20.26\\($+1.3\%$)} & \shortstack{20.89 / 20.24\\($+4.5\%$ / $+1.2\%$)} & \shortstack{20.89 / 19.88\\($+4.5\%$ / $-0.6\%$)} & \shortstack{17.29 / 17.71\\($-13.5\%$ / $-11.4\%$)} \\[3pt]
小孔（5） & \shortstack{5.017\\($+0.3\%$)} & \shortstack{5.08 / 5.013\\($+1.7\%$ / $+0.3\%$)} & \shortstack{4.16 / 4.19\\($-16.8\%$ / $-16.2\%$)} & \shortstack{3.01 / 2.96\\($-39.8\%$ / $-40.9\%$)} \\[3pt]
喉道（0.05） & \shortstack{0.0503\\($+0.6\%$)} & \shortstack{0.0508 / 0.0505\\($+1.5\%$ / $+1.1\%$)} & 未分辨 & 未分辨 \\
\bottomrule
\end{tabular}
\end{table}

\begin{itemize}
\item 方式 1 与方式 2 的 $T_{2\mathrm{lm}}$ 几乎一致：孤立孔本身可以被准确反算；方式 2 的主峰取分箱中心（每箱约 $2.4\%$），因此略偏大。
\item 方式 3 中小孔分量只占 $5.9\%$，且与占 $94\%$ 的大孔 $T_2$ 仅差约 $2.8$ 倍，正则化平滑把小峰推向短 $T_2$，孔径偏小约 $16\%$。
  这是反演方法造成的偏差，不是物理效应。
\item 方式 4 在反演偏差之外，还叠加了孔间扩散耦合与喉道壁的弛豫“汇”：直接拟合连通链中小孔区域的磁化得 $T_2=\SI{312}{ms}$（孤立时 \SI{430}{ms}），
  反算孔径偏小约 $40\%$。
\item 喉道仅占孔面积的 $0.22\%$，在方式 3、4 中未被分辨。
\end{itemize}

\section{适用边界}

\begin{itemize}
\item 以上为孤立孔的解；连通体系中孔隙之间通过喉道的扩散交换以及喉道壁的“弛豫汇”会使小孔 $T_2$ 变短（本案例数值模拟中由 \SI{430}{ms} 降至约 \SI{312}{ms}），不能直接套用孤立孔公式换算尺寸。
\item 二维圆盘对应无限长圆柱，$S/V=2/R$；三维球形孔为 $S/V=3/R$，二者换算系数不同。
\item 快扩散判据 $\kappa=\rho_2a/D\ll1$；本案例 $\kappa\le0.05$，故快扩散近似误差 $<1\%$。
\end{itemize}

\begin{thebibliography}{8}
\setlength{\itemsep}{2pt}

\bibitem[Brownstein \& Tarr(1979)]{BrownsteinTarr1979}
Brownstein, K. R., \& Tarr, C. E. (1979). Importance of classical diffusion in NMR studies of water in biological cells.
\emph{Physical Review A}, \emph{19}(6), 2446--2453. \url{https://doi.org/10.1103/PhysRevA.19.2446}

\bibitem[Carslaw \& Jaeger(1959)]{CarslawJaeger1959}
Carslaw, H. S., \& Jaeger, J. C. (1959). \emph{Conduction of heat in solids} (2nd ed.). Oxford, UK: Clarendon Press.

\bibitem[Coates et~al.(1999)]{Coates1999}
Coates, G. R., Xiao, L., \& Prammer, M. G. (1999). \emph{NMR logging: Principles and applications}.
Houston, TX: Halliburton Energy Services.

\bibitem[Grebenkov(2007)]{Grebenkov2007}
Grebenkov, D. S. (2007). NMR survey of reflected Brownian motion. \emph{Reviews of Modern Physics}, \emph{79}(3), 1077--1137.
\url{https://doi.org/10.1103/RevModPhys.79.1077}

\bibitem[Kleinberg et~al.(1994)]{Kleinberg1994}
Kleinberg, R. L., Kenyon, W. E., \& Mitra, P. P. (1994). Mechanism of NMR relaxation of fluids in rock.
\emph{Journal of Magnetic Resonance, Series A}, \emph{108}(2), 206--214. \url{https://doi.org/10.1006/jmra.1994.1112}

\bibitem[Senturia \& Robinson(1970)]{Senturia1970}
Senturia, S. D., \& Robinson, J. D. (1970). Nuclear spin-lattice relaxation of liquids confined in porous solids.
\emph{Society of Petroleum Engineers Journal}, \emph{10}(3), 237--244. \url{https://doi.org/10.2118/2870-PA}

\bibitem[Torrey(1956)]{Torrey1956}
Torrey, H. C. (1956). Bloch equations with diffusion terms. \emph{Physical Review}, \emph{104}(3), 563--565.
\url{https://doi.org/10.1103/PhysRev.104.563}

\bibitem[Zhou \& Zhang(2025)]{Zhou2025}
Zhou, \& Zhang (2025). How does pore structure affect the NMR relaxation in unsaturated porous media.
\emph{Water Resources Research}. （作者全名、卷期与 DOI 请按原文核对；$D$、$\TB$、$\rho_2$ 的参数来源。）
\end{thebibliography}

\end{document}
